\documentclass[11pt]{letter}

\usepackage[margin=0.5in]{geometry}
\usepackage{hyperref}

\begin{document}

\begin{letter}{}

\opening{Dear Hiring Team,}

I am writing to apply for the Machine Learning Research Engineer position, role number 200670886-4102, with Apple's Product Operations machine learning team in Waterloo. With a Master of Science in Computer Science from the University of Calgary and applied experience in Python, PyTorch, multimodal data processing, model evaluation, and scalable research software, I am excited by the opportunity to contribute to machine learning research for manufacturing data and production-oriented ML capabilities.

My graduate research focused on developing machine learning and data-management systems for large geospatial and Earth observation datasets. In my digital elevation model super-resolution project, I built Python and PyTorch workflows for data collection, preprocessing, model training, and evaluation. I trained deep ResNet models for 4x terrain reconstruction and compared model quality using RMSE, slope, TRI, TPI, and wavelet-based metrics. This work required careful experimentation, debugging, and evaluation of whether model performance generalized beyond convenient random splits.

I also developed C++ and Python software for scalable processing, storage, and retrieval of multimodal satellite imagery through my research work with BigGeo, Mitacs, and the University of Calgary. I automated SAR and optical imagery preprocessing, built DGGS-based encoding pipelines, designed tensor-based storage for multiresolution retrieval, and improved processing performance with C++ multithreading. These projects are different from factory data, but they gave me strong experience with messy, high-volume, multimodal data and the engineering discipline needed to make research methods usable.

The Apple role's focus on anomaly detection, automated machine learning, time-series, graph, image, and tabular data is especially interesting to me. My background gives me a strong base in deep learning, scientific data workflows, evaluation, and research implementation. I am also comfortable collaborating across research and engineering contexts, from publishing a first-author paper to mentoring deep learning students and building tools that help others inspect, compare, and use model outputs.

I would be glad to bring my PyTorch, Python, data-processing, research software, and technical communication experience to Apple's Product Operations machine learning team while continuing to grow in manufacturing ML, anomaly detection, AutoML, and production translation of successful research ideas.

Thank you for considering my application. I would welcome the opportunity to discuss how my research engineering background can contribute to Apple's machine learning work in Waterloo.

\closing{Sincerely,

Amir Mirzai Golpayegani

\href{mailto:amirgolpaa24@gmail.com}{amirgolpaa24@gmail.com}}

\end{letter}

\end{document}
